% GENERATED FILE. Edit canonical lessons, then run npm run generate:books.
% canonical-source-hash: fnv1a64:eed51c08737eca6a
% canonical-lessons: ZH-C35-gongzuo, ZH-C35-lao, ZH-C35-laoshi, ZH-C35-xiansheng, ZH-C35-zhu

\chapter{Who Does What}
\label{ch:zh-who-does-what}
\hlchaptermodality{35}

\begin{chapteropening}
\textbf{By the end of this chapter:} \emph{I can say what people do and where somebody lives.}

Complete the last lesson of chapter 35: I can say what people do and where somebody lives.
\end{chapteropening}

\section[gōngzuò]{gōngzuò — to work; job}

\label{lesson:ZH-C35-gongzuo}



\begin{quote}
Before the new one: say the Mandarin for to know someone, then the Mandarin for friend.
\end{quote}

\subsection*{You'll want to know: gōngzuò}
\emph{gōngzuò} — \textquotedblleft{}to work; job\textquotedblright{}.

\begin{quote}
\emph{wǒ jiějie gōngzuò, wǒ shàngxué} — My sister works; I go to school.
\end{quote}

Its characters come later; for now it is a word for the ear.

\begin{cousinweb}
\emph{gōngzuò} is built from gōng 'craft' and zuò 'to do', both inherited from Middle Chinese, \emph{kuwng} and \emph{tsak} (Guangyun, in Baxter's transcription); the dictionary glosses the first 'office' and 'skill', the second 'to do'.
\end{cousinweb}

\begin{grammarlens}[title={What you've built}]
One more word about people.
\end{grammarlens}

\subsection*{Your turn}
\begin{itemize}
\raggedright
  \item \emph{Say it:} \emph{gōngzuò}
  \item \emph{Say it:} \emph{gōngzuò}, once more
  \item \emph{Say it:} \emph{wǒ jiějie gōngzuò, wǒ shàngxué} — My sister works; I go to school.
  \item \emph{Recall:} say the Mandarin for to know someone, then the Mandarin for friend, then say \emph{gōngzuò} again
\end{itemize}

\begin{tcolorbox}[breakable,colback=teal!4,colframe=teal!35!black,title={Before you move on}]
What does gōngzuò mean? (\textquotedblleft{}To work; job\textquotedblright{}.) Say it once more.
\end{tcolorbox}
\section[lǎo]{lǎo — old}

\label{lesson:ZH-C35-lao}



\begin{quote}
Before the new one: say the Mandarin for friend, then the Mandarin for to work.
\end{quote}

\subsection*{You'll want to know: lǎo}
\emph{lǎo} — \textquotedblleft{}old\textquotedblright{}.

\begin{quote}
\emph{tā shì wǒ de lǎo tóngxué} — He's an old classmate of mine.
\end{quote}

Its characters come later; for now it is a word for the ear.

\begin{cousinweb}
\emph{lǎo} is inherited from Middle Chinese \emph{lawX} (Guangyun, in Baxter's transcription), glossed 'aged'.
\end{cousinweb}

\begin{grammarlens}[title={What you've built}]
One more word about people.
\end{grammarlens}

\subsection*{Your turn}
\begin{itemize}
\raggedright
  \item \emph{Say it:} \emph{lǎo}
  \item \emph{Say it:} \emph{lǎo}, once more
  \item \emph{Say it:} \emph{tā shì wǒ de lǎo tóngxué} — He's an old classmate of mine.
  \item \emph{Recall:} say the Mandarin for friend, then the Mandarin for to work, then say \emph{lǎo} again
\end{itemize}

\begin{tcolorbox}[breakable,colback=teal!4,colframe=teal!35!black,title={Before you move on}]
What does lǎo mean? (\textquotedblleft{}Old\textquotedblright{}.) Say it once more.
\end{tcolorbox}
\section[lǎoshī]{lǎoshī — teacher}

\label{lesson:ZH-C35-laoshi}



\begin{quote}
Before the new one: say the Mandarin for to work, then the Mandarin for old.
\end{quote}

\subsection*{You'll want to know: lǎoshī}
\emph{lǎoshī} — \textquotedblleft{}teacher\textquotedblright{}.

\begin{quote}
\emph{wǒ māma shì lǎoshī} — My mum is a teacher.
\end{quote}

Its characters come later; for now it is a word for the ear.

\begin{cousinweb}
\emph{lǎoshī} is built from lǎo 'old, respected' and shī 'teacher', both inherited from Middle Chinese: \emph{lawX}, glossed 'aged', and \emph{srij}, glossed 'a model to follow' (Guangyun, in Baxter's transcription).
\end{cousinweb}

\begin{grammarlens}[title={What you've built}]
One more word about people.
\end{grammarlens}

\subsection*{Your turn}
\begin{itemize}
\raggedright
  \item \emph{Say it:} \emph{lǎoshī}
  \item \emph{Say it:} \emph{lǎoshī}, once more
  \item \emph{Say it:} \emph{wǒ māma shì lǎoshī} — My mum is a teacher.
  \item \emph{Recall:} say the Mandarin for to work, then the Mandarin for old, then say \emph{lǎoshī} again
\end{itemize}

\begin{tcolorbox}[breakable,colback=teal!4,colframe=teal!35!black,title={Before you move on}]
What does lǎoshī mean? (\textquotedblleft{}Teacher\textquotedblright{}.) Say it once more.
\end{tcolorbox}
\section[xiānsheng]{xiānsheng — Mr; sir; husband}

\label{lesson:ZH-C35-xiansheng}



\begin{quote}
Before the new one: say the Mandarin for old, then the Mandarin for teacher.
\end{quote}

\subsection*{You'll want to know: xiānsheng}
\emph{xiānsheng} — \textquotedblleft{}Mr; sir; husband\textquotedblright{}.

\begin{quote}
\emph{tā xiānsheng shì lǎoshī} — Her husband is a teacher.
\end{quote}

Its characters come later; for now it is a word for the ear.

\begin{cousinweb}
\emph{xiānsheng} is built from xiān 'before' and \zh{生} 'to be born, to grow', both inherited from Middle Chinese, \emph{sen} and \emph{srjaeng} (Guangyun, in Baxter's transcription). 'The one born first' is the usual explanation of how it came to mean an older man to respect, and then Mr or husband; the path is not settled.
\end{cousinweb}

\begin{grammarlens}[title={What you've built}]
One more word about people.
\end{grammarlens}

\subsection*{Your turn}
\begin{itemize}
\raggedright
  \item \emph{Say it:} \emph{xiānsheng}
  \item \emph{Say it:} \emph{xiānsheng}, once more
  \item \emph{Say it:} \emph{tā xiānsheng shì lǎoshī} — Her husband is a teacher.
  \item \emph{Recall:} say the Mandarin for old, then the Mandarin for teacher, then say \emph{xiānsheng} again
\end{itemize}

\begin{tcolorbox}[breakable,colback=teal!4,colframe=teal!35!black,title={Before you move on}]
What does xiānsheng mean? (\textquotedblleft{}Mr; sir; husband\textquotedblright{}.) Say it once more.
\end{tcolorbox}
\section[zhù]{zhù — to live somewhere; to stay}

\label{lesson:ZH-C35-zhu}



\begin{quote}
Before the new one: say the Mandarin for teacher, then the Mandarin for sir.
\end{quote}

\subsection*{You'll want to know: zhù}
\emph{zhù} — \textquotedblleft{}to live somewhere; to stay\textquotedblright{}.

\begin{quote}
\emph{tā zhù wǒ jiā} — He's staying at our place.
\end{quote}

Its characters come later; for now it is a word for the ear.

\begin{cousinweb}
\emph{zhù} is inherited from Middle Chinese \emph{drjuH} (Guangyun, in Baxter's transcription), glossed 'to stop, to stay'. Staying somewhere becoming living there is a likely step, but it is not settled.
\end{cousinweb}

\begin{grammarlens}[title={What you've built}]
One more word about people.
\end{grammarlens}

\subsection*{Your turn}
\begin{itemize}
\raggedright
  \item \emph{Say it:} \emph{zhù}
  \item \emph{Say it:} \emph{zhù}, once more
  \item \emph{Say it:} \emph{tā zhù wǒ jiā} — He's staying at our place.
  \item \emph{Recall:} say the Mandarin for teacher, then the Mandarin for sir, then say \emph{zhù} again
\end{itemize}

\begin{tcolorbox}[breakable,colback=teal!4,colframe=teal!35!black,title={Before you move on}]
What does zhù mean? (\textquotedblleft{}To live somewhere; to stay\textquotedblright{}.) Say it once more.
\end{tcolorbox}
